\section{A big-step semantics for \QGA}
\label{sec:semantics}

\providecommand{\tv}[1]{\ensuremath{\mathsf{#1}}}
\providecommand{\bs}[1]{\ensuremath{\Downarrow_{#1}}}

\S\ref{sec:reduction} is one route to a self-consistency result.  This section
sketches the other, which is to give \QGA{} a semantics of its own and repeat the
argument of \S\ref{sec:proof} with \QGA{} in place of \BGA.  Nothing here is
formalized yet.  The purpose is to fix the shape of the semantics before the
encoding work begins, because one clause in it determines the architecture of
everything that follows.

\subsection{What the proof needs}

The consistency argument of \S\ref{sec:proof} has three moving parts.
Evaluation is deterministic, the proof checker is sound with respect to
evaluation, and \fn{False} evaluates to false.  Consistency is then immediate:
a proof of \fn{False} would make \fn{False} evaluate to true, and determinism
forbids that.

Two constraints follow, and they are the constraints that matter.  The semantics
has to be expressible as a \QGA{} recursive definition, since it will live inside
\QGA{} rather than above it.  And the satisfaction predicate has to be decided at
each fixed index, because the soundness proof is an induction whose motive
carries satisfaction of the hypotheses in an antecedent position, and $\longrightarrow$
introduction demands that the antecedent be grounded.  This second constraint is
the whole reason \S\ref{sec:fuel} indexes evaluation by fuel, and it is the
constraint that decides the design below.

\subsection{Terms}

Term evaluation is the semantics of \S\ref{sec:fuel} unchanged.  A judgment
$\langle t, A\rangle \bs{k} v$ says that the term coded by $t$, under the
assignment $A$, evaluates to the value $v$ within $k$ steps.  It is realized as
a \QGA{} function \fn{eval\_fuel} returning $\Suc{v}$ on success and $0$ on
timeout, which makes it total, so equality of results is decidable and the
groundedness obligations discharge themselves.  Divergence is the absence of any
$k$, and the relation $\langle t, A\rangle \Downarrow v$ abbreviates
$\exists k.\ \langle t, A\rangle \bs{k} v$.  Fuel is monotone, so a successful
evaluation stays successful when more is supplied.

The definition list is a parameter of the semantics.  A constant introduced by
$f\;\bar{x} \rec \fn{body}$ is evaluated by looking its body up and continuing,
exactly as \fn{dfns} is consulted in \S\ref{sec:bga}.

\subsection{Formulas}
\label{sec:semantics:formulas}

Formula evaluation is two-valued with gaps.  A formula either evaluates to
\tv{tt}, or to \tv{ff}, or to neither, and the third case is what groundedness
means.  Writing $\langle \varphi, A \rangle \Downarrow \tv{tt}$ and $\langle
\varphi, A \rangle \Downarrow \tv{ff}$ for the two judgments, the clauses for the
primitives are the strong three-valued ones.

\begin{center}
\begin{tabular}{@{}ll@{}}
\toprule
$\langle s = t, A\rangle \Downarrow \tv{tt}$ &
  if $\langle s, A\rangle \Downarrow n$ and $\langle t, A\rangle \Downarrow n$\\
$\langle s = t, A\rangle \Downarrow \tv{ff}$ &
  if $\langle s, A\rangle \Downarrow m$, $\langle t, A\rangle \Downarrow n$ and
  $m \neq n$\\
\midrule
$\langle \neg\varphi, A\rangle \Downarrow \tv{tt}$ &
  if $\langle \varphi, A\rangle \Downarrow \tv{ff}$\\
$\langle \neg\varphi, A\rangle \Downarrow \tv{ff}$ &
  if $\langle \varphi, A\rangle \Downarrow \tv{tt}$\\
\midrule
$\langle \varphi \vee \psi, A\rangle \Downarrow \tv{tt}$ &
  if $\langle \varphi, A\rangle \Downarrow \tv{tt}$, or if
  $\langle \psi, A\rangle \Downarrow \tv{tt}$\\
$\langle \varphi \vee \psi, A\rangle \Downarrow \tv{ff}$ &
  if $\langle \varphi, A\rangle \Downarrow \tv{ff}$ and
  $\langle \psi, A\rangle \Downarrow \tv{ff}$\\
\midrule
$\langle \ite{c}{\varphi}{\psi}, A\rangle \Downarrow v$ &
  if $\langle c, A\rangle \Downarrow \tv{tt}$ and
  $\langle \varphi, A\rangle \Downarrow v$\\
$\langle \ite{c}{\varphi}{\psi}, A\rangle \Downarrow v$ &
  if $\langle c, A\rangle \Downarrow \tv{ff}$ and
  $\langle \psi, A\rangle \Downarrow v$\\
\bottomrule
\end{tabular}
\end{center}

\noindent
Each clause is positive, so the whole definition is a monotone induction.  An
equation is true when both sides converge to the same value and false when they
converge to different ones, which is where habeas quid comes from: a divergent
side gives neither clause, so the equation is ungrounded, and \isa{a N} is
grounded exactly when $a$ converges.  Disjunction is true as soon as one disjunct
is, whatever the other does, which is what makes $\vee$I1 and $\vee$I2 sound
while excluded middle is not.  Nothing evaluates to \tv{tt} and \tv{ff} together,
which is determinism and is proved by induction on the two derivations.

\subsection{The quantifier clause}
\label{sec:semantics:quant}

Six of the eight clauses above are bounded searches, so they are $\Sigma_1$ and a
\QGA{} recursive definition can perform them.  The quantifiers are where a choice
has to be made, and the two candidate readings differ in exactly the property
that \S\ref{sec:semantics} identified as decisive.

The first reading is instance-wise.

\begin{center}
\begin{tabular}{@{}ll@{}}
$\langle \forall x.\,\varphi, A\rangle \Downarrow \tv{tt}$ &
  if $\langle \varphi, A[x \mapsto n]\rangle \Downarrow \tv{tt}$ for every $n$\\
$\langle \forall x.\,\varphi, A\rangle \Downarrow \tv{ff}$ &
  if $\langle \varphi, A[x \mapsto n]\rangle \Downarrow \tv{ff}$ for some $n$\\
\end{tabular}
\end{center}

\noindent
This is the reading a classical reader expects, and it makes soundness of
$\forall$I a substitution lemma and nothing more.  \QGA{} is first order, its
terms denote numbers or nothing, and it has no term-level binder, so assignments
suffice and no strengthening of the induction motive is required.  Determinism at
the quantifier is equally easy, since one instance cannot be both true and false.

Its defect is that the true clause quantifies over all $n$.  The predicate is not
$\Sigma_1$, so no \QGA{} definition decides it, and the internal soundness proof
would have satisfaction sitting undecided in an antecedent.  That is precisely
the obstruction \S\ref{sec:proof} had to work around for \BGA, and there the
obstruction was removable because the undecidability came from an unbounded
search for fuel.  Here it comes from a universal quantification over the
numerals, and fuel does not touch it.

The second reading grounds a universal in the system's own proof search.

\begin{center}
\begin{tabular}{@{}ll@{}}
$\langle \forall x.\,\varphi, A\rangle \Downarrow \tv{tt}$ &
  if some valid \QGA{} proof derives $x \N \seq \varphi$\\
$\langle \forall x.\,\varphi, A\rangle \Downarrow \tv{ff}$ &
  if some valid \QGA{} proof derives $\seq \neg\varphi[x \mapsto \bar{n}]$ for
  some numeral $\bar{n}$\\
\end{tabular}
\end{center}

\noindent
Both clauses are searches over proof codes, so the predicate stays $\Sigma_1$ and
a fuelled version of it is total and decided at each index, which is what the
soundness induction needs.  Both mention derivability positively, so the
simultaneous definition of the checker and the semantics remains a monotone
induction and no Tarskian circularity arises.  The clauses for $\exists$ are the
duals, with the true clause a search for a witness and the false clause a search
for a proof of the schematic negation.

This reading also settles the two axioms of \S\ref{sec:quant}.
$\neg\forall$I is the false clause read from right to left, since a derivation of
$\neg(Q\,a)$ together with $a \N$ is the certificate the clause asks for, and
$\neg\forall$E is the same clause read from left to right.  Under the
instance-wise reading they are equally sound.  Under either one they are
justified rather than postulated, which is the reason they were added.

We propose the second reading.  The cost is the architecture it forces.  The
soundness case for $\forall$I has to produce a certificate with no residual
hypotheses, while the induction hypothesis supplies only that the premise
judgment is semantically true, so the hypotheses have to be converted back into
derivations and discharged by \textsc{cut}.  That step is an internal
completeness lemma, and it has to be established before soundness rather than
after.  Reflective grounded arithmetic is organized this way for the same reason,
and it is the main risk in the plan.

\subsection{Obligations}

Assuming the second reading, a self-consistency proof for \QGA{} needs the
following, in this order.

\begin{enumerate}
\item An arithmetization of \QGA's own syntax, judgments, proofs and proof
      checker, following \S\ref{sec:bga}.  This is needed to state consistency at
      all, and the reflective clauses need it as well.
\item The fuelled evaluation of terms and formulas, with monotonicity in the
      index, and the decidedness lemma for satisfaction at a fixed index.
\item Internal completeness: a formula that evaluates to \tv{tt} is derivable.
\item Soundness: every rule of \S\ref{sec:qga} preserves truth.  The propositional
      cases read off the clauses directly, the equality cases need a congruence
      lemma for evaluation, \textsc{Ind} needs induction on the value of its
      subject, and \textsc{defE} and \textsc{defI} need that a definition and its
      body evaluate alike.
\item Determinism, which under the reflective reading rests on soundness, because
      a colliding pair of certificates has to be instantiated at a common numeral
      before the contradiction appears.
\item $\langle \fn{False}, A\rangle \Downarrow \tv{ff}$, which is one evaluation.
\end{enumerate}

\noindent
Consistency follows from the last three.  Two of these items deserve a warning.
Item 3 is the one \S\ref{sec:semantics:quant} identified as the cost of the
design.  Item 4 has a trap of its own at the substitution rule: the obvious
motive, truth under all satisfying assignments, is not inductive when
substitution can replace a variable by a term that changes what the hypotheses
say.  For a first-order system the standard motive is expected to suffice, and
confirming that is part of the work.

\subsection{What the semantics does not settle}

Grounded truth under either reading is only recursively enumerable, so the
diagonal argument produces an ungrounded sentence rather than a true and
unprovable one, and a system that is consistent, effective and complete for its
own semantics is not excluded.  That is why the arrangement is available at all.

What remains open is whether the resulting proof deserves the name
self-consistency.  The semantics interprets \QGA's $\forall$ by a search over
\QGA{} proofs, so the argument is carried out in the same system it is about, and
a Kripkean partial truth predicate is what makes that admissible.  A sceptical
reader is entitled to ask for the adequacy statement relating the arithmetized
syntax of item 1 to the Pure-level formulas of \S\ref{sec:levels}, and the
present development does not supply it.  \S\ref{sec:reduction} avoids the
question by a different route.
